\documentclass[11pt,a4paper]{article}\usepackage{layout}\begin{document}
\begin{canvas}
\node[anchor=north west,inner sep=0pt] at (current page.north west){\includegraphics[width=210mm,height=297mm]{.build/assets/cover-base.pdf}};
\Box{22}{182.4}{166}{\centering\sffamily\bfseries\fontsize{11.8}{15}\selectfont\color{white}精要版\hspace{4mm}{\digits\fontsize{9.6}{12}\selectfont ESSENTIAL EDITION}}
\Box{22}{191.0}{166}{\centering\sffamily\fontsize{9.2}{12}\selectfont\color{white}从普惠到质量，从参与到影响}
\node[anchor=west,inner sep=0pt] at \Pos{27}{279.9}{\includegraphics[width=29mm,trim=83bp 84bp 84bp 83bp,clip]{assets/xlab-ai.png}};
\node[anchor=west,inner sep=0pt] at \Pos{78}{279.9}{\includegraphics[width=56mm,trim=17bp 70bp 18bp 69bp,clip]{assets/siri.png}};
\node[anchor=west,inner sep=0pt] at \Pos{155}{279.9}{\includegraphics[width=30mm,trim=0bp 0bp 0bp 0bp,clip]{assets/ccosi.png}};
\end{canvas}
\begin{canvas}\Header{奇点智能研究院 · 定位与使命}{teal}\Footer{03}
\Box{15}{28}{180}{\fontsize{21.5}{29}\selectfont\color{navy}\sffamily \bfseries 奇点智能研究院}
\Box{15}{44}{180}{\fontsize{11.4}{16}\selectfont\color{teal}\sffamily \bfseries Singularity Intelligence Research Institute}
\Box{15}{57}{180}{\fontsize{12.5}{18}\selectfont\color{ink}\sffamily \bfseries 人工智能产业的“范式孵化器”：连接前沿研究、产业洞察与实践转化。}
\Box{15}{74}{180}{\includegraphics[width=180mm]{.build/assets/institute-research-system.pdf}}
\Box{15}{178}{180}{\fontsize{8.1}{12}\selectfont\color{muted}\sffamily 图 I　研究体系：问题、证据与实践相互反馈。根据机构介绍整理。}
\Box{15}{192}{86}{\fontsize{11.4}{17}\selectfont\color{blue}\sffamily \bfseries 由技术变化出发}
\Box{15}{201}{86}{\fontsize{10.5}{17}\selectfont\color{ink}奇点智能研究院由专业开发者社区 CSDN 与 IT 技术咨询平台 Boolan 联合发起，定位于创新型研究机构、科技智库与战略咨询平台。研究院联结高校、科研机构、产业企业与开发者社区，开展跨学科研究。}
\Box{109}{192}{86}{\fontsize{11.4}{17}\selectfont\color{teal}\sffamily \bfseries 以产业与社会价值为方向}
\Box{109}{201}{86}{\fontsize{10.5}{17}\selectfont\color{ink}研究院关注人工智能与开源创新如何共同改变技术底座、产业生态和社会治理。其使命是理解范式转换的逻辑，打通技术突破与产业应用的路径，使研究能够回应真实问题。}
\Rect{15}{253}{180}{23}{paper}
\Box{20}{257}{170}{\fontsize{10.8}{15}\selectfont\color{teal}\sffamily \bfseries 愿景}
\Box{40}{256.8}{150}{\fontsize{10.6}{16.8}\selectfont\color{ink}孕育思想、发展方法，推动智能技术成为普惠、可信、可持续的全球公共品。}
\end{canvas}
\begin{canvas}\Header{奇点智能研究院 · 五大范式}{teal}\Footer{04}
\Box{15}{28}{180}{\fontsize{21.5}{29}\selectfont\color{navy}\sffamily \bfseries 五大范式转换}
\Box{15}{44}{180}{\fontsize{11.4}{16}\selectfont\color{teal}\sffamily \bfseries 从技术底座与研发方式，延伸到开放协作和共同治理}
\Box{15}{57}{180}{\fontsize{10.5}{17}\selectfont\color{ink}“范式”指组织技术与工作的基本方式。研究院用五个相互关联的观察方向，理解大模型带来的变化，并辨认变化真正落地所需的条件。}
\Box{15}{77}{180}{\includegraphics[width=180mm]{.build/assets/institute-five-paradigms.pdf}}
\Box{15}{214}{180}{\fontsize{8.1}{12}\selectfont\color{muted}\sffamily 图 II　五大范式的观察坐标。根据机构介绍整理；箭头表示研究方向的扩展。}
\Box{15}{226}{86}{\fontsize{11.4}{17}\selectfont\color{blue}\sffamily \bfseries 变化发生在不同层次}
\Box{15}{235}{86}{\fontsize{10.5}{17}\selectfont\color{ink}计算关注基础架构，开发与交互关注人怎样完成任务；协作与治理则把问题扩展到资源如何共享、责任怎样落实。技术进步需要与组织方式相互适配。}
\Box{109}{226}{86}{\fontsize{11.4}{17}\selectfont\color{teal}\sffamily \bfseries 开放贯穿五个方向}
\Box{109}{235}{86}{\fontsize{10.5}{17}\selectfont\color{ink}开放资源提供学习与复用的基础，社区连接不同主体，透明的规则支持信任。开放能否转为公共能力，还取决于使用门槛、持续维护与实际服务结果。}
\end{canvas}
\begin{canvas}\Header{奇点智能研究院 · 研究布局}{teal}\Footer{05}
\Box{15}{28}{180}{\fontsize{21.5}{29}\selectfont\color{navy}\sffamily \bfseries 六大研究方向}
\Box{15}{44}{180}{\fontsize{11.4}{16}\selectfont\color{teal}\sffamily \bfseries 把技术能力、开发者生态与可信治理放在同一研究框架中}
\Box{15}{57}{180}{\fontsize{10.5}{17}\selectfont\color{ink}六个方向分别回应“系统怎样运行、团队怎样协作、组织怎样转型”以及“如何观察生态、改善参与、建立信任”。它们通过数据、方法和实际场景相互连接。}
\Box{15}{77}{180}{\includegraphics[width=180mm]{.build/assets/institute-six-directions.pdf}}
\Box{15}{207}{180}{\fontsize{8.1}{12}\selectfont\color{muted}\sffamily 图 III　六大研究方向的关联结构。根据奇点智能研究院机构介绍整理。}
\Box{15}{221}{86}{\fontsize{11.4}{17}\selectfont\color{blue}\sffamily \bfseries 让生态具备“可开发性”}
\Box{15}{230}{86}{\fontsize{10.5}{17}\selectfont\color{ink}基础软件的能力，需要通过工具、文档、社区与协作机制被开发者使用。研究院将生态可开发性与开发者体验联读，关注参与门槛、贡献路径和持续成长。}
\Box{109}{221}{86}{\fontsize{11.4}{17}\selectfont\color{teal}\sffamily \bfseries 让智能转型具有治理基础}
\Box{109}{230}{86}{\fontsize{10.5}{17}\selectfont\color{ink}研发和业务流程的变化，同时需要许可合规、供应链安全及责任机制。AI 物料清单（AIBOM）用于记录 AI 系统的组成与依赖，为追踪风险和持续维护提供基础。}
\end{canvas}
\begin{canvas}\Header{奇点智能研究院 · 研究与实践}{teal}\Footer{06}
\Box{15}{28}{180}{\fontsize{21.5}{29}\selectfont\color{navy}\sffamily \bfseries 从研究问题走向可检验的实践}
\Box{15}{44}{180}{\fontsize{11.4}{16}\selectfont\color{teal}\sffamily \bfseries 以数据理解现状，以方法组织行动，以结果修正判断}
\Box{15}{57}{180}{\fontsize{10.5}{17}\selectfont\color{ink}研究院面向政府、企业与开发者开展研究与咨询。不同对象的问题各有侧重，但都需要把证据、方法和实践连接起来，才能将前沿判断转为可以讨论和检验的行动。}
\Box{15}{77}{180}{\includegraphics[width=180mm]{.build/assets/institute-practice-path.pdf}}
\Box{15}{206}{180}{\fontsize{8.1}{12}\selectfont\color{muted}\sffamily 图 IV　研究进入实践的路径。根据机构介绍与本报告的公共能力框架整理。}
\Box{15}{219}{86}{\fontsize{11.4}{17}\selectfont\color{blue}\sffamily \bfseries AI 原生研发：形成改进的方法}
\Box{15}{228}{86}{\fontsize{10.5}{17}\selectfont\color{ink}AI 原生软件研发成熟度模型（AISMM）白皮书已在研究院官网公开。它围绕研发全生命周期、人机协作与组织流程展开，为团队识别能力差距、选择改进路径提供框架。}
\Box{109}{219}{86}{\fontsize{11.4}{17}\selectfont\color{teal}\sffamily \bfseries 全球开源洞察：形成决策的证据}
\Box{109}{228}{86}{\fontsize{10.5}{17}\selectfont\color{ink}本报告将生态地图、主体榜单、技术全景与模型网络联读，沿着参与、供给、复用、治理与公共价值展开，讨论开放资源怎样转化为可持续公共能力。}
\Box{15}{271.7}{180}{\fontsize{7.8}{11}\selectfont\color{muted}\sffamily 资料依据：奇点智能研究院机构介绍；《AI 原生软件研发成熟度模型》白皮书公开介绍。}
\end{canvas}
\end{document}